\chapter{Software Model}
\label{ch:sw}

The directory \file{sw\_model} holds a software model of the units in C, with C++
and Python interfaces, test programs and example applications. The model computes
the result of each operation, with no timing. For the formats 16/5, 32/8 and 64/11
(\sig{DATA\_W}/\sig{EXP\_W}) it is bit-accurate to the RTL, as far as the check of
Section~\ref{sec:sw_check} covers. Other formats are accepted but not verified.
\sig{SPECIAL\_CASES}~=~0 is not modelled. The commands are in
Section~\ref{sec:make_sw}, and Table~\ref{tab:sw_files} lists the contents.

\begin{twotable}{Software Model Contents}{tab:sw_files}{Path}{Content}{4.6cm}
\file{c/ieee754\_model.h}, \file{c/ieee754\_model.c} & C model. \\ \hline
\file{cpp/ieee754.h} & C++ types \sig{Ieee754} and \sig{Ieee754Mac}. \\ \hline
\file{python/ieee754.py} & Python wrapper of the C model. \\ \hline
\file{python/floatVSieee754} & Tool that compares two output files. \\ \hline
\file{tests/ieee754\_tests.cpp} & Self-checking test program. \\ \hline
\file{tests/check.py}, \file{tests/model\_run.c}, \file{tests/rtl\_tb.v} & Comparison with TestFloat and the RTL. \\ \hline
\file{examples} & Example applications. \\
\end{twotable}

\section{C Model}
\label{sec:sw_c}

Operands and results are raw encodings in the low \sig{data\_w} bits of a
\sig{uint64\_t}. Rounding is to nearest, ties to even, and there are no exception
flags. Every NaN result has sign 0 and only the top fraction bit set.
Table~\ref{tab:sw_c} lists the functions. The accepted formats have
$6 \le \mathtt{data\_w}-\mathtt{exp\_w} \le 53$, $2 \le \mathtt{exp\_w} \le 11$ and
$\mathtt{data\_w} \le 64$.

\begin{twotable}{C Model Functions}{tab:sw_c}{Function}{Operation}{6.2cm}
\sig{ieee754\_format\_ok(data\_w, exp\_w)} & Returns 1 if the format is accepted. \\ \hline
\sig{ieee754\_add}, \sig{ieee754\_sub}, \sig{ieee754\_mul}, \sig{ieee754\_div} & Operation on \sig{(a, b, data\_w, exp\_w)}. \\ \hline
\sig{ieee754\_from\_double(value, data\_w, exp\_w)} & Conversion from binary64, rounding to nearest even. \\ \hline
\sig{ieee754\_to\_double(bits, data\_w, exp\_w)} & Conversion to binary64. \\ \hline
\sig{ieee754\_mac\_new(data\_w, exp\_w, lguard\_w, rguard\_w)} & Creates a clocked model of the \sig{mac} pipeline. Returns \sig{NULL} for an unaccepted format or if $\mathtt{lguard\_w}+\mathtt{rguard\_w}+2(\mathtt{data\_w}-\mathtt{exp\_w})+5 > 128$. \\ \hline
\sig{ieee754\_mac\_push(mac, a, b)} & Applies one product of the current stream. \\ \hline
\sig{ieee754\_mac\_result(mac)} & Ends the stream and returns its sum; $+0$ for an empty stream. \\ \hline
\sig{ieee754\_mac\_free(mac)} & Frees the model. \\
\end{twotable}

\section{C++ Interface}
\label{sec:sw_cpp}

\file{cpp/ieee754.h} defines two classes, for the format given by the macros
\sig{IEEE754\_DATA\_W} (default 32) and \sig{IEEE754\_EXP\_W} (default 8).
\sig{Ieee754} holds one encoding, converts from and to the arithmetic types, and
has arithmetic and comparison operators. Addition, subtraction, multiplication and
division run on the model, and math functions go through binary64.
\sig{Ieee754Mac} is the streaming \sig{fmac}, with the methods \sig{push} and
\sig{result}.

\section{Python Interface}
\label{sec:sw_py}

\file{python/ieee754.py} wraps the shared library \file{libieee754.so}, which the
software Makefile builds in \file{sw\_model/build}. The environment variable
\sig{IEEE754\_LIB} can give another path. The class \sig{Format} has the methods
\sig{add}, \sig{sub}, \sig{mul}, \sig{div}, \sig{from\_double}, \sig{to\_double}
and \sig{mac}, and the class \sig{Mac} is the streaming \sig{fmac}.

\file{python/floatVSieee754} compares a reference file in binary64 with a target
file in \sig{float} or in the model format. For each differing line it writes the
error and the symmetric relative error. It ends with the maximum and mean of both,
and the numbers of differing and total lines.

\section{Self-Checking Tests}
\label{sec:sw_tests}

\file{tests/ieee754\_tests.cpp} needs no RTL tools. It tests the conversions, the
arithmetic operators, the comparators, memory copies and the streaming MAC. It
compares 200,000 conversions and 200,000 random operand pairs per operation with
the \sig{float} arithmetic of the host. It prints PASS or FAIL.

\section{Comparison With the RTL}
\label{sec:sw_check}

\file{tests/check.py} compares the model with TestFloat and with the RTL, simulated
with Icarus Verilog. A comparison is bit for bit, except that any NaN equals any
NaN in the TestFloat comparison. It covers the following, and
Table~\ref{tab:sw_check} gives the sizes:

\begin{itemize}
\item the C model against all TestFloat level 1 vectors of add, sub, mul and div;
\item the C model against the \sig{fpu} RTL, on a random mix of the four operations;
\item the C model against the \sig{fmac} RTL, on product streams with several
  \sig{LGUARD\_W} and \sig{RGUARD\_W} values;
\item the Python wrapper against the C model, on the same mixes and streams.
\end{itemize}

\begin{threetable}{Model Comparison With the RTL}{tab:sw_check}{Format}{\sig{fpu} operations}{\sig{fmac} streams, for each (\sig{LGUARD\_W}, \sig{RGUARD\_W})}{2cm}{3.4cm}
32/8 & 100,000 & 6,000, for (0,0), (4,0), (0,4), (3,5), (8,8) \\ \hline
16/5 & 20,000 & 2,000, for (0,0), (2,3) \\ \hline
64/11 & 20,000 & 2,000, for (0,0), (2,3) \\
\end{threetable}

The check runs the \sig{fpu} and \sig{fmac} RTL. \sig{fadd}, \sig{fmul} and
\sig{fdiv} use the same \sig{add}, \sig{mul} and \sig{div} datapaths, and are not
simulated by the check.

\section{Example Applications}
\label{sec:sw_examples}

The directory \file{examples} has the six applications of Table~\ref{tab:sw_apps}.
Each one computes its result in \sig{double}, \sig{float} and model arithmetic and
writes each to a file in \file{sw\_model/build/out}. \sig{make run} then compares
the \sig{float} and model files with the \sig{double} file.

\begin{twotable}{Example Applications}{tab:sw_apps}{\sig{CPROG}}{Application}{3.4cm}
\sig{dp} & Dot product. \\ \hline
\sig{newton-raph} & Newton-Raphson. \\ \hline
\sig{fft} & Fast Fourier transform. \\ \hline
\sig{k-means} & K-means clustering. \\ \hline
\sig{black-scholes} & Black-Scholes. \\ \hline
\sig{pink-noise} & Pink noise. \\
\end{twotable}
